\documentclass[10pt, a4paper]{article}
\usepackage{preamble}

\title{Topology II}
\author{Luke Phillips}
\date{August 2025}

\begin{document}

\maketitle

\newpage

\tableofcontents

\newpage

\section{Metric Spaces}

\begin{definition}[Metric space]
    A metric space is a pair $(X, d)$ where $X$ is a set and $d : X \times X \to [0, \infty)$ is a function such that,
    for all $x, y, z \in X$:

    \begin{enumerate}[label = (\alph*)]
        \item
        \[
        d(x, y) = 0 \iff x = y.
        \]

        \item
        (symmetry)
        \[
        d(x, y) = d(y, x)\quad\forall x, y \in X.
        \]

        \item
        (triangle inequality)
        \[
        d(x, y) \leq d(x, z) + d(z, y)\quad\forall x, y, z \in X.
        \]
    \end{enumerate}
\end{definition}

We call $d$ a metric on $X$.

\begin{example}
    Let $p \in [1, \infty)$,
    and $n \in \N$.
    Then $\R ^ n$ has a metric
    \[
    d_p(x, y) = (|x_1 - y_1| ^ p + \dotsi + |x_n - y_n| ^ p) ^ {1 / p} = \left(\sum_{i = 1}^{n}|x_i - y_i| ^ p\right) ^ {\frac{1}{p}}.
    \]
\end{example}

\begin{example}
    Let $X$ be any set.
    Then
    \[
    d(x, y) = \begin{cases}
        0 &\text{if } x = y \\
        1 &\text{if } x \neq y
    \end{cases}
    \]
    is a metric,
    called the discrete metric on $X$.
\end{example}

\begin{example}
    Let $p \in [1, \infty)$,
    and $n \in \N$.
    Then $\R ^ n$ has a metric
    \[
    d_p(x, y) = (|x_1 - y_1| ^ p + \dotsi + |x_n - y_n| ^ p) ^ {1 / p} = \left(\sum_{i = 1}^{n}|x_i - y_i| ^ p\right) ^ {\frac{1}{p}}.
    \]
    Is a metric space.

    \begin{solution}
        Symmetry is easily derived from the properties of modulus.

        \begin{align*}
            d_p(x, y) = 0 &\iff \left(\sum_{i = 1}^{n}|x_i - y_i| ^ p\right) ^ {\frac{1}{p}} = 0 \\
            &\iff \sum_{i = 1}^{n}|x_i - y_i| ^ p = 0 \\
            \intertext{Since $0 \leq |x_i - y_i|$ for all $i \leq n$}
            &\iff \forall i \leq n, |x_i - y_i| = 0 \\
            &\iff \forall i \leq n, x_i = y_i \\
            &\iff x = y.
        \end{align*}

        For the triangle inequality start with the following:
        set $a_k = |x_k - z_k|$,
        $b_k = |z_k - y_k|$.
        Clearly $0 \leq a_k, b_k$.
        Then
        \begin{align*}
            d_p(x, y) &= \left(\sum_{i = 1}^{n}|x_i - y_i| ^ p\right) ^ {\frac{1}{p}} \\
            &= \left(\sum_{i = 1}^{n}|(x_i - z_i) + (z_i - y_i)| ^ p\right) ^ {\frac{1}{p}} \\
            &\leq \left(\sum_{i = 1}^{n}(|x_i - z_i| + |z_i - y_i|) ^ p\right) ^ {\frac{1}{p}} \\
            &= \left(\sum_{i = 1}^{n}(a_i + b_i) ^ p\right) ^ {\frac{1}{p}} \\
            &\leq \left(\sum_{i = 1}^{n}a_i ^ p\right) ^ {\frac{1}{p}} + \left(\sum_{i = 1}^{n}b_i ^ p\right) ^ {\frac{1}{p}} &\text{(by Minkowski's inequality)} \\
            &= d_p(x, z) + d_p(z, y).
        \end{align*}
    \end{solution}
\end{example}

\begin{definition}
    Let $(X, d)$ be a metric space,
    $r > 0$ and $x \in X$.
    Then
    \[
    B(x; r) = \{y \in X\mid d(x, y) < r\}
    \]
    is called the open ball of radius $r$ around $x$,
    and
    \[
    D(x; r) = \{y \in X\mid d(x, y) \leq r\}
    \]
    is called the closed ball of radius $r$ around $x$.
\end{definition}

\subsection{Open and closed sets}

\begin{definition}
    Let $X$ be a metric space.
    A subset $U \subset X$ is called open,
    if for every $x \in U$ there is $\epsilon > 0$ such that $B(x; \epsilon) \subset U$.
    A subset $A \subset X$ is called closed,
    if $X \setminus A$ is open.
\end{definition}

\begin{example}
    Let $\R$ have the usual $|x - y|$ metric.
    Then,
    considered as sets:
    
    \begin{enumerate}[label = \arabic*.]
        \item Any single point $\{x\} \in \R$ is not open.

        \item Any single point $\{x\} \in \R$ is closed.
    \end{enumerate}

    \begin{proof}
        \begin{enumerate}[label = \arabic*.]
            \item
            We have $\{x\} \in \R$ so for $\{x\}$ to be not open we need
            \begin{align*}
                \{x\} \subset \R, \neg(\forall x \in \{x\}, \exists \varepsilon > 0, B(x; \varepsilon) \subset \{x\}) &\iff \{x\} \subset \R, \forall \varepsilon > 0, B(x; \varepsilon) \not\subset \{x\} \\
                &\iff \forall \varepsilon > 0, \{y \in X \mid |x - y| < \varepsilon\} \not\subset \{x\}. \\
                &\iff \forall\varepsilon > 0, \exists y \in X, |x - y| < \varepsilon \wedge y \neq x.
            \end{align*}
            So let $\varepsilon > 0$,
            then we need to find a $y \in X$ such that $|x - y| < \varepsilon$ such that $y \notin \{x\}$,
            we can simply choose $y = x + \varepsilon / 2$ as
            \[
            |x - y| = |x - (x + \varepsilon / 2)| = |\varepsilon / 2| < \varepsilon
            \]
            hence $y \in B(x, \varepsilon)$,
            but $y \notin \{x\}$.
            So $\{x\}$ is not open.

            \item
            For $\{x\} \in \R$ to be closed we need for $\R \setminus \{x\}$ to be open and hence
            \[
            \forall y \in \R \setminus \{x\}, \exists \varepsilon > 0, B(y; \varepsilon) \subset \R \setminus \{x\} \iff \forall y \neq x, \exists \varepsilon > 0, |x - y| < \varepsilon
            \]
            hence we can choose $\varepsilon = |x - y| > 0$,
            which means we have $x \notin B(y; \varepsilon)$,
            $y \in B(y; \varepsilon) \subseteq \R \setminus \{x\}$.
            So $\R \setminus \{x\}$ is open,
            so $\{x\}$ is closed.
        \end{enumerate}
    \end{proof}
\end{example}

\begin{definition}
    Suppose $X$ is a metric space and $x \in N \subseteq X$.
    We call $N$ a neighbourhood of $x$ if there exists an open set $V \subseteq X$ with $x \in V \subseteq N$.
\end{definition}

From the previous two definitions we can get the following:
\begin{corollary}
    For $x$ in a metric space $X$,
    $N \subseteq X$ is a neighbourhood of $x$ if and only if there is some $\varepsilon > 0$ such that $x \in B(x; \varepsilon) \subseteq N$.

    A set $U \subseteq X$ is open if and only if $U$ is a neighbourhood of each $x \in U$.
\end{corollary}

\begin{proposition}
    Open balls are open,
    and closed balls are closed.

    \begin{proof}
        Let $B(x; r)$ be an open ball,
        with $x \in X$ and $r > 0$,
        suppose $y \in B(x; r)$.
        We need to find an open ball round $y$ within $B(x; r)$.
        So let $s = d(x, y)$.
        Then we know $s < r$ and we can set $\varepsilon = r - s > 0$.
        Now we claim that $B(t; \varepsilon) \subset B(x; r)$.
        To show this,
        consider any $z \in B(y; \varepsilon)$.
        Then
        \begin{align*}
            d(z, x) &\leq d(z, y) + d(y, x) \\
            &< \varepsilon + s \\
            &= (r - s) + s = r.
        \end{align*}
        So $z \in B(x; r)$,
        and we have $y \in B(y; \varepsilon) \subset B(x; r)$ as required.

        Now let $D(x; r)$ be a closed ball.
        To show it is closed,
        we must show $X \setminus D(x; r)$ is open.

        If $X \setminus D(x; r) = \emptyset$ there is nothing to show.
        So we consider some $y \in X \setminus D(x; r)$.
        This has $s = d(x, y) > r$.
        Then setting $\varepsilon = s - r > 0$,
        we claim that $B(y; \varepsilon) \subset X \setminus D(x; r)$.
        Let $z \in B(y; \varepsilon)$.
        Then
        \begin{align*}
            d(z, x) &\geq d(x, y) - d(y, z) \\
            &> s - \varepsilon \\
            &= s - (s - r) = r.
        \end{align*}
        So $z \notin D(x; r)$,
        and we have $y \in B(y; \varepsilon) \subset X \setminus D(x; r)$ as required.
        This means $X \setminus D(x; r)$ is open,
        so $D(x; r)$ is closed.
    \end{proof}
\end{proposition}

\begin{lemma}
    Let $X$ be a metric space.

    \begin{enumerate}[label = (\alph*)]
        \item $X$ and $\emptyset$ are both open and closed.

        \item An arbitrary union of open sets is open.

        \item A finite intersection of open sets is open.

        \item A finite union of closed sets is closed.

        \item An arbitrary intersection of closed sets is closed.
    \end{enumerate}

    \begin{proof}
        \begin{enumerate}[label = (\alph*)]
            \item
            For $x \in X$,
            we have $x \in B(x; 1) \subseteq X$,
            as
            \[
            x \in \{y \in X \mid |x - y| < 1\} \subseteq X.
            \]
            Also,
            there is no point $x$ in $\emptyset$.
            So $\forall x \in \emptyset, \exists \varepsilon > 0, \{y \in X \mid |x - y| < \varepsilon\} \subset X$,
            since there is no $x$ we can trivially see that this definition is satisfied by merely taking $\varepsilon = |y|$ for any $y \in X$.

            Both are closed as $X$ is closed if $X \setminus X = \emptyset$ is open which we have just shown,
            $\emptyset$ is closed if $X \setminus \emptyset = X$ is open which we have just shown.

            \item
            Let $X_i \subset X$ be open for all $i \in I$
            ($I$ is some interval).
            Then if $x \in \bigcup_{i \in I}X_i$ we have $x \in X_{r}$ for some $r \in I$.
            So $\exists \varepsilon > 0, a \in B(x; \varepsilon) \subseteq X_{r} \subseteq \bigcup_{i \in I}X_i$.
            Hence $\bigcup_{i \in I}X_i$ is open.

            \item
            Let $U_i \subset X$ be open for $i = 1, \dotsc, n$.
            Then if $x \in \bigcap_{i = 1}^{n}U_i$ we have $x \in U_i$ for $i = 1, \dotsc, n$.
            So $\forall i, \exists \varepsilon_i > 0, x \in B(x; \varepsilon_i) \subseteq U_i$.
            Let $\varepsilon = \min\{\varepsilon_1, \dotsc, \varepsilon_n\} > 0$.
            Then for $i = 1, \dotsc, n$ we have $x \in B(x; \varepsilon) \subseteq B(x; \varepsilon_i) \subseteq U_i$.
            Hence $x \in B(x; \varepsilon) \subseteq \bigcap_{i = 1}^{n}U_i$ as required.

            \item
            Let $V_i \subset X$ be closed for $i = 1, \dotsc, n$.
            Then by definition $X \setminus V_i \subseteq X$ is open for $i = 1, \dotsc, n$,
            so by the previous
            (c)
            we know $X \setminus \left(\bigcup_{i = 1}^{n}V_i\right) = \bigcap_{i = 1}^{n}(X \setminus V_i)$ is open,
            so $\bigcup_{i = 1}^{n}V_i$ is closed as required.
        \end{enumerate}
    \end{proof}
\end{lemma}

\subsection{Continuity}

\begin{definition}
    Let $(X, d)$ be a metric space.

    \begin{enumerate}[label = (\alph*)]
        \item A sequence in $X$ is a function $a : \N = \{0, 1, \dotsc\} \to X$.
        We write $(a_n)_{n \in \N}$ or simply $(a_n)$ for the sequence,
        where $a_n = a(n)$.

        \item Let $a \in X$.
        Then the sequence $(a_n)_{n \in \N}$ converges to $a$,
        if for all $\varepsilon > 0$ there is a $N \in \N$ with $d(a, a_n) < \varepsilon$ for all $n \geq N$.
        We write $\lim_{n \to \infty}a_n = a$ or $a_n \to a$.
        
        In symbols this is:
        \[
        \forall \varepsilon > 0, \exists N \in \N, \forall n \geq N, d(a, a_n) < \varepsilon.
        \]
    \end{enumerate}
\end{definition}

We would say "converges in $d$ to $a$" written as $a_n \xrightarrow{d} a$.

















\end{document}